% LSPV-3 REPORT::BOTTLENECK1::2026-10-03T03:38:00-07:00
\documentclass[11pt]{article}
\usepackage[margin=.88in]{geometry}
\usepackage{amsmath,amssymb}
\begin{document}
\section*{LSPV-3 BOTTLENECK1: exact value and ceiling of route optimization}
\textbf{STATUS: OPEN.} P versus NP remains open. This report proves deterministic accounting identities only. It proves no SAT lower bound, no arbitrary-algorithm acquisition theorem, and no barrier-crossing hinge.

\section*{CORPUS REVIEWED}
The current repository corpus was surveyed before substantive work. Governance read includes AGENTS.md, CODEX.md, the pnp-research skill, and the corpus-reconciliation note. Root main.tex and verified MAIN 0058 were inspected, together with current task surfaces including TASK-004 and TASK-005. Materially relevant LSPV work was reconciled: LSPV-1 structural PAIRDEG1/FULLPAIR6 lineage; LSPV-2 EXCESS/EXCESSDUAL1/TURANEX1 lineage; LSPV-3 FACTOR1, QLEN1, RAMCOMP1, ENVELOPE1, RAMDOM1, CROSSDOM1, CERTREG1 and TOTALENV1; and LSPV-4 TRACECLOSE1, LEGALFLAT1, RESPSEL1 and GLOBALGAME1. HEC, CRS-PPC and CCSPC were surveyed. The immediate LSPV-2 predecessor surface remains verified. Repository search found no exact preserved failed peer scientific body requiring a recovery write; historical missing bodies were not reconstructed.

Primary literature was refreshed. Viola's alternating-time simulations and oracle obstruction remain relativizing. Van Melkebeek--Mocelin Sdroievski remains the targeted-HSG/AM white-box source. Akshay et al. prove polynomial Skolem/Herbrand extraction from proofs in their declared QBF proof system; proof acquisition is not free here. Current DQBF dependency-scheme work remains proof-system/model specific.

\section*{DUPLICATION / CONTRADICTION AUDIT}
TOTALENV1 already proves
\[
T_{\rm total}(N)\le N^{\max\{G,E(K)\}+o(1)}
\]
for common-charge exponent \(G\) and route envelope
\[
E(K)=\min\{D(K),R(K),F(K)\}.
\]
The corpus states qualitatively that route improvements are invisible when \(G\) dominates. BOTTLENECK1 does not duplicate the envelope. It quantifies the exact maximum total-exponent benefit obtainable from any route-only improvement and identifies when route optimization is exhausted. No contradiction with GLOBALGAME1 was found: GLOBALGAME1 concerns acquisition/continuation viability, whereas BOTTLENECK1 conditions on declared executable routes.

\section*{BOTTLENECK1 exact theorem}
Fix the SAME\_REALIZATION common-charge exponent \(G\ge0\). Let the current cheapest route exponent be \(E\ge0\), so the total exponent is
\[
H=\max\{G,E\}.
\]
Suppose a valid reconstruction improvement changes only route-specific work, replacing \(E\) by \(E'\) with \(0\le E'\le E\), while all genuinely common obligations and their exponent \(G\) are unchanged. Then the total exponent improvement is exactly
\[
\boxed{\Delta_H=\max\{G,E\}-\max\{G,E'\}.}
\]
Consequently
\[
\boxed{0\le\Delta_H\le \max\{0,E-G\}.}
\]
The upper bound is attained whenever \(E'\le G\). Therefore the maximum possible total-exponent gain from \emph{all future route-only optimization}, even an idealized zero-exponent route, is
\[
\boxed{\Delta_{\rm route}^{\max}=\max\{0,E-G\}.}
\]
This is \textbf{PROVED/EXACT} as polynomial-exponent accounting under the declared common/route partition.

Three regimes follow exactly.
\[
E<G:\quad H=G,\quad \Delta_H=0\ \text{for every }E'\le E.
\]
\[
E=G:\quad H=G,\quad \Delta_H=0\ \text{for every }E'\le E.
\]
\[
E>G:\quad H=E,\quad
\Delta_H=
\begin{cases}
E-E',&E'\ge G,\\
E-G,&E'\le G.
\end{cases}
\]
Thus once a route has been improved to exponent at most \(G\), further route optimization is asymptotically invisible at polynomial-exponent scale unless a common obligation is also reduced.

\section*{Application to ENVELOPE1}
With
\[
D(K)=\max\{a,q+bK\},\quad
R(K)=\ell+c_0K^2,\quad
F(K)=\max\{f,c_FK\},
\]
and
\[
E(K)=\min\{D(K),R(K),F(K)\},
\]
the exact remaining value of route research at fixed certified \(K\) is bounded by
\[
\boxed{\Delta_{\rm route}^{\max}(K)=\max\{0,E(K)-G\}.}
\]
This is a prioritization theorem, not a lower bound. If \(E(K)\le G\), source effort should move from route crossovers to the common obligation attaining \(G\), unless route work is needed for a different proof obligation such as SAME\_REALIZATION acquisition.

\section*{Quantifier and reconstruction audit}
The load-bearing predicate remains
\[
\exists\sigma\;\forall u\;\exists z\;\forall y\;R(x,Q,\sigma,u,z,y),
\]
with one shared size-normalized \(Q\), Arthur randomness, exact SAT YES/NO conditions, and model-matched compilation. BOTTLENECK1 performs no quantifier elimination. FACTOR1 remains one exact-decider layer for a fixed answer vector. Generic quantifier commutation remains \textbf{FALSIFIED/RETIRED}. Generic canonical binding remains unavailable.

GLOBALGAME1 adds a useful guard: if a globally coupled response route depends on a continuation predicate, that route is not executable merely because a local response exists. Such acquisition/viability costs must be charged before the route enters \(E(K)\).

\section*{Full charge and TransportCert}
The time-level ledger remains
\[
\tau_Q+\tau_{\rm norm}+\tau_{\rm census}+\tau_{\rm amp}+\tau_{\rm idx}
+\tau_{\rm common\,compile}+\min\{\tau_D,\tau_R,\tau_F\}
+\tau_{\rm write}+\tau_{\rm transport}.
\]
A cost belongs in \(G\) only if every declared implementation must pay that SAME\_REALIZATION obligation. Route-specific evaluator acquisition, continuation acquisition, compilation, or indexing stays in the relevant route. TransportCert remains accounting only, not a lower-bound theorem.

\section*{PQ and cross-domain classification}
Construction-specific collision extraction or special-soundness binding remains \textbf{ANALOGY/CANDIDATE} until an unconditional source-matched theorem exists. Generic cryptographic binding and random-oracle arguments remain \textbf{FAILED TRANSFER} for unconditional P versus NP. QBF proof-system Skolem extraction is a \textbf{KNOWN THEOREM} in its declared proof model but only an \textbf{ANALOGY/CANDIDATE} for arbitrary SAT computation because proof acquisition is uncharged otherwise. No PQ technique reviewed changes BOTTLENECK1's accounting identity or supplies TC4.

\section*{Barrier audit}
BOTTLENECK1 relativizes and survives algebraic-oracle augmentation. A semantic oracle can supply route outputs or common objects. The theorem does not force an arbitrary SAT decider to materialize \(Q\), a selector, evaluator, proof, continuation strategy, or transcript. SAME\_REALIZATION acquisition, TC4 arbitrary-algorithm escape, nonrelativization, and nonalgebrization remain \textbf{OPEN}.

\section*{WHAT CHANGED RELATIVE TO ROOT MAIN}
Root main.tex records ENVELOPE1, RAMDOM1 and CROSSDOM1 but its LSPV-3 region is stale relative to verified CERTREG1 and TOTALENV1. BOTTLENECK1 adds the exact ceiling
\[
\Delta_{\rm route}^{\max}=\max\{0,E-G\}
\]
on total-exponent improvement obtainable by changing only route-specific work. It therefore supplies a stopping rule for route optimization once \(E\le G\).

\section*{ADOPTED / REJECTED / RETIRED PRIOR COMPONENTS}
\textbf{Adopted:} FACTOR1, QLEN1, RAMCOMP1, ENVELOPE1, RAMDOM1, CROSSDOM1, CERTREG1, TOTALENV1, EXPINV1, explicit acquisition charges, GLOBALGAME1 as a route-viability guard, and the common/route partition.

\textbf{Rejected as free:} generic canonical binding, evaluator or continuation acquisition, proof acquisition, qualitative PH collapse as a numerical compiler, exponent-certificate inflation, and moving route-specific costs into the common ledger.

\textbf{FALSIFIED/RETIRED:} generic quantifier commutation and SGLAUD1 constant-query amplification remain retired.

\section*{STALE OR OVERSTATED CLAIMS FOUND}
Any claim that a route improvement necessarily improves total reconstruction time is overstated. TOTALENV1 already requires comparison with \(G\); BOTTLENECK1 quantifies the exact maximum benefit. No verified corpus artifact establishes an unrestricted SAT lower bound from these accounting envelopes.

\section*{Next target}
Source-transcribe the actual vMS/SGL SAME\_REALIZATION tuple
\[
(a,q,b,\ell,c_0,f,c_F,K)
\]
and the common exponents
\[
(g_Q,g_{\rm norm},g_{\rm census},g_{\rm amp},g_{\rm idx},
g_{\rm common\,compile},g_{\rm write},g_{\rm transport}).
\]
Compute \(E(K)\), \(G\), and \(\Delta_{\rm route}^{\max}\). If \(E(K)\le G\), stop abstract route optimization and attack the common obligation attaining \(G\). If \(E(K)>G\), the remaining route-improvement budget is exactly \(E(K)-G\).

\section*{References}
E. Viola, On Probabilistic Time versus Alternating Time, ECCC TR05-137, 2005.

D. van Melkebeek and N. Mocelin Sdroievski, Instance-Wise Hardness and Refutation versus Derandomization for Arthur-Merlin Protocols, Computational Complexity 34:15, 2025.

S. Akshay, O. Beyersdorff, S. Chakraborty, L. Kasche, M. Mahajan, and L. N. Spachmann, Proof Systems for QBF Synthesis: Extracting Skolem and Herbrand Functions, SAT 2026, LIPIcs 377, Article 3.

L. Chew and T. Peitl, Strong (D)QBF Dependency Schemes via Pure Paths with Applications to Proof Checking, SAT 2026, LIPIcs 377, Article 11.
\end{document}
% END_OF_VERIFIED_REPORT::LSPV-3::BOTTLENECK1